% AI-produced by GPT-6 Astra (OpenAI, ChatGPT Pro mode).
\subsection{Realizing the four-dimensional component obstruction}

\begin{saresult}{A.140}{Proved under stated hypotheses}{Realizing the component obstruction and evaluating it for simply connected four-manifolds}
Let $B$ be a connected closed oriented real PL four-manifold, and
retain $I_B,J_B,\lambda_B$ and $c_P$ from A.137. Put
$\pi=\pi_1(B)$ and let $\mathbb Z_{(2)}$ denote the integers
localized at the prime $2$. Define
\[
 \begin{split}
 \delta_* &: H_2(B;\mathbb Z/2)\longrightarrow
                         H_1(B;\mathbb Z_{(2)}),\\
 L_B&=\mathrm{PD}_B^{-1}(\ker\delta_*)
                      \subseteq H^2(B;\mathbb Z/2),
 \end{split}                                             \tag{A.140.1}
\]
where $\delta_*$ is the Bockstein for
$0\to\mathbb Z_{(2)}\xrightarrow{2}\mathbb Z_{(2)}
\to\mathbb Z/2\to0$, and $\mathrm{PD}_B$ is mod-$2$
Poincar\'e duality. Suppose that the assembly summand of
Orson--Powell--Randal-Williams, Definition~5.8, is zero:
\[
 I_2:H_2(\pi;\mathbb Z_{(2)})
       \longrightarrow L^q_6(\mathbb Z[\pi])_{(2)},
 \qquad I_2=0.                                          \tag{A.140.2}
\]
Then
\[
 \operatorname{im}\lambda_B
       \supseteq (L_B+J_B)/J_B.                          \tag{A.140.3}
\]
In particular, if $H_1(B;\mathbb Z_{(2)})$ has no $2$-torsion,
$\lambda_B$ is surjective and A.137 becomes the isomorphism
\[
 I_B/\operatorname{im}(\pi_0c_P)
          \cong H^2(B;\mathbb Z/2)/J_B.                 \tag{A.140.4}
\]
This does not evaluate the endpoint-image condition defining $I_B$
or the subgroup $J_B$ in general.

If $B$ is simply connected, both of those restrictions disappear
on the TOP concordance component group, and
\[
       \pi_0P_{\mathrm{TOP}}(B)=1,
       \qquad J_B=H^2(B;\mathbb Z/2).                   \tag{A.140.5}
\]
There is the stronger ordinary-loop assertion. If $\gamma$ is a
based TOP loop of automorphisms of $B$, let
$T_\gamma(x,t)=(\gamma_t(x),t)$ and define
\[
 \eta_B:\pi_1 A_{\mathrm{TOP}}(B)\longrightarrow
                  H^2(B;\mathbb Z/2),
 \qquad \eta_B[\gamma]=\kappa(T_\gamma).                \tag{A.140.6}
\]
For simply connected $B$ this homomorphism is surjective and
vanishes on the image of $\pi_1A_{\mathrm{PL}}(B)$.
Consequently the ordinary automorphism comparison
$A_{\mathrm{PL}}(B)\to A_{\mathrm{TOP}}(B)$ is not surjective
on $\pi_1$ whenever $H^2(B;\mathbb Z/2)\ne0$. In particular
this holds for the finite PL manifold $S^2\times S^2$.
This last statement concerns automorphisms of the base; it is
not an additional conclusion about the homotopy groups of
$P_{\mathrm{PL}}(B)\to P_{\mathrm{TOP}}(B)$.

No computation of $\pi_0P_{\mathrm{PL}}(B)$, no smooth isotopy
conclusion in dimension four, and no full higher TOP concordance
computation is asserted.
\end{saresult}
\begin{proof}
Choose a compatible smoothing of the given PL four-manifold.
The existence of such a smoothing is stated explicitly in the
last paragraph of Kirby--Siebenmann, Essay~I, p.~39
\cite{KirbySiebenmann1977}. Only existence for this single
manifold is used. No PL isotopy is converted into a smooth isotopy.
Let $\sigma$ denote the resulting product smooth structure on
$N=B\times I$, and let $\alpha$ be its underlying standard PL
structure. We use common product collars as in A.136.

\textbf{Identifying the two obstruction conventions.}
Suppose first that $F:N\to N$ is a TOP concordance fixing both
boundary components. Its A.137 invariant is
\[
 \kappa(F)=[F_*\alpha]-[\alpha]
       \in H^3(N,\partial N;\mathbb Z/2)
       \cong H^2(B;\mathbb Z/2).                       \tag{A.140.7}
\]
Set $W=N\times I$ and
$C=(N\times\{0\})\cup(\partial N\times I)$.
The boundary smoothing used by Orson--Powell--Randal-Williams
is standard on $C$ and is $F^*\sigma$ on $N\times\{1\}$.
Their original definition, pp.~2 and 7--9, assigns to this boundary
smoothing its relative universal Kirby--Siebenmann obstruction in
$H^4(W,\partial W;\mathbb Z/2)$ and then applies Poincar\'e
duality to obtain $\mathrm{KS}_{\mathrm{OPR}}(F)$ in
$H_2(B;\mathbb Z/2)$
\cite[\S3.1, proof of Lemma~2.2, p.~9]{OrsonPowellRandalWilliams2026}.

Here is the exact suspension identification. The standard lift on
$W$ trivializes the torsor of PL lifts. Changing only the top
boundary lift by
\[
 d=[F^*\alpha]-[\alpha]
       \in H^3(N,\partial N;\mathbb Z/2)
       \cong H^3(\partial W,C;\mathbb Z/2)
\]
changes the relative universal obstruction by the triple connecting
map
\[
 \partial_{!W}:
 H^3(\partial W,C;\mathbb Z/2)
       \xrightarrow{\cong}H^4(W,\partial W;\mathbb Z/2).
                                                               \tag{A.140.8}
\]
Indeed the universal square of Kirby--Siebenmann, Essay~V,
\S5.6, pp.~252--253, models a PL lift as a nullhomotopy of the
universal map to $K(\mathbb Z/2,4)$. Two boundary nullhomotopies
which agree on $C$ differ by a map
$(\partial W,C)\to(K(\mathbb Z/2,3),*)$; its obstruction to
extending over $W$ is precisely the relative connecting class
in (A.140.8). This describes the actual relative classes, not only
their zero sets. The same square is obtained by applying the
map from smooth lifts to PL lifts to the obstruction used in
the cited definition of $\mathrm{KS}_{\mathrm{OPR}}$.
The map in (A.140.8) is an isomorphism because the inclusion
$C\hookrightarrow W$ is a homotopy equivalence: both spaces
deformation retract to $N\times\{0\}$. It is exactly the
interval suspension under the displayed identifications.

Since $F^*\alpha=(F^{-1})_*\alpha$, the class $d$ is
$\kappa(F^{-1})$. A.137 proves that the action of a bottom-fixed
concordance on $H^3(N,\partial N;\mathbb Z/2)$ is trivial, and
that $\kappa$ is additive on boundary-PL concordances. Hence
$\kappa(F^{-1})=-\kappa(F)=\kappa(F)$, with these coefficients.
Finally the two interval suspensions and mod-$2$ Poincar\'e
duality give
\[
       \mathrm{KS}_{\mathrm{OPR}}(F)
                 =\mathrm{PD}_B\bigl(\kappa(F)\bigr).
                                                               \tag{A.140.9}
\]
The equality is unchanged by the common-collar normalization:
that normalization is an isotopy relative to the boundary and
therefore preserves both relative lift classes. Thus no boundary
regularity has been added silently.

\textbf{What the block theorem realizes.}
Theorem~2.3(D) of Orson--Powell--Randal-Williams, printed p.~6,
under exactly (A.140.2), says that the image of
\[
 \pi_1\widetilde{\operatorname{Homeo}}^+(B)
       \xrightarrow{a}Q(B)
       \xrightarrow{\mathrm{KS}_{\mathrm{OPR}}}
                       H_2(B;\mathbb Z/2)
\]
contains $\ker\delta_*$. The tilde is essential: these are
\emph{block} homeomorphisms, as defined in their \S2.1,
pp.~4--5. They are not ordinary loops of homeomorphisms.
There is a source-notation discrepancy to make explicit:
Theorem~2.3(D) prints the codomain as $L^s_6$, whereas the
referenced Definition~5.8, p.~21, and its proof on p.~22 specify
the quadratic assembly summand with codomain $L^q_6$.
Equation (A.140.2) uses that actual definition and proof, with
the quadratic convention introduced in their \S5.2.1.

We spell out why the theorem supplies a literal both-end-fixed
TOP concordance for every class $z\in\ker\delta_*$.
Choose $q=a(u)$ with $\mathrm{KS}_{\mathrm{OPR}}(q)=z$,
and represent $q$ by a TOP concordance $F$ from the identity to
a diffeomorphism $f$. Exactness of their displayed sequence (1),
p.~6, says that the class of $f$ is trivial in
$\pi_0\widetilde{\operatorname{Diffeo}}^+(B)$.
The explicit interpretation in \S2.1 is smooth pseudo-isotopy,
so choose a smooth pseudo-isotopy $G$ from $f$ to the identity.
The smooth concordance
\[
 G'=G\circ(f^{-1}\times1)
\]
has bottom identity and top $f^{-1}$. Consequently
$F_0=G'\circ F$ has both boundary maps literally equal to the
identity. The original Lemma~3.2 and the subsequent definition
on p.~9 make $\mathrm{KS}_{\mathrm{OPR}}$ additive on $Q(B)$;
it is zero on the smooth concordance $G'$. Therefore
\[
 \mathrm{KS}_{\mathrm{OPR}}(F_0)=z,
 \qquad \kappa(F_0)=\mathrm{PD}_B^{-1}(z).              \tag{A.140.10}
\]
We have used smoothness of $G'$ only to make its smoothing
obstruction zero. We have not called $G'$ literally PL in fixed
coordinates; the PL boundary condition needed for (A.140.9)
holds for $F_0$ because its two endpoints are the identity.

The ordinary free-top component of $F_0$ belongs to $I_B$,
since its top is identity. By A.137 its image under $\lambda_B$
is $\mathrm{PD}_B^{-1}(z)+J_B$. This proves (A.140.3).
If $H_1(B;\mathbb Z_{(2)})$ has no $2$-torsion, exactness of
the coefficient sequence gives $\delta_*=0$. Hence $L_B=H^2(B;
\mathbb Z/2)$ and $\lambda_B$ is onto. Its kernel was computed
in A.137, proving (A.140.4). None of these steps places a block
loop directly in $J_B$.

\textbf{The simply connected case and ordinary loops.}
Suppose now that $B$ is simply connected. It is orientable;
$H_2(\pi;\mathbb Z_{(2)})=0$ makes (A.140.2) automatic, and
$H_1(B;\mathbb Z_{(2)})=0$ makes $\delta_*=0$.
The corrected theorem of Gabai--Gay--Hartman--Krushkal--Powell,
Theorem~1.5, states that every TOP pseudo-isotopy of a compact
simply connected four-manifold is TOP-isotopic to the identity
relative to the bottom and the vertical boundary
\cite{GabaiGayHartmanKrushkalPowell2026}. This is a statement
about the whole ordinary pseudo-isotopy, not just its endpoint.
For our closed $B$ it says exactly that $P_{\mathrm{TOP}}(B)$
is path connected. In particular each top map is TOP-isotopic
to the identity, so the endpoint predicate defining $I_B$ is
automatic. Since $I_B$ is then the trivial group, the surjectivity
in (A.140.4) already implies $J_B=H^2(B;\mathbb Z/2)$.

For the stronger assertion, take any prescribed
$v\in H^2(B;\mathbb Z/2)$ and use (A.140.10) to choose a
both-end-fixed $F_0$ with $\kappa(F_0)=v$.
Choose an ordinary path $H_s$ from the identity to $F_0$ by
the cited theorem, and put $\gamma_s=e(H_s)$.
Because both endpoint maps of $F_0$ are identity, $\gamma$ is
an ordinary based loop in $A_{\mathrm{TOP}}(B)$.
There are explicit cylinder paths
\[
 T_s(x,t)=(\gamma_{st}(x),t),
 \qquad Q_s=T_s^{-1}H_s.                               \tag{A.140.11}
\]
Each $Q_s$ fixes both boundary components. It starts at identity
and ends at $T_\gamma^{-1}F_0$. Apply the common-collar
normalization of A.136 to the whole finite parameter family.
The relative class $\kappa(Q_s)$ is then constant, and is zero.
Additivity in A.137 gives
\[
                 \kappa(T_\gamma)=\kappa(F_0)=v.       \tag{A.140.12}
\]
This proves surjectivity using genuine ordinary loops, obtained
only after the simply connected theorem has been applied.

Homotopies of based loops give relative-boundary families of their
traces, so (A.140.6) is well defined. Pointwise multiplication of
loops agrees with multiplication of the traces; A.137 additivity
therefore proves that $\eta_B$ is a homomorphism. More precisely
it is the composite
\[
 \pi_1A_{\mathrm{TOP}}(B)\longrightarrow\pi_0F_A
       \xrightarrow{\pi_0\beta_B}H^2(B;\mathbb Z/2),  \tag{A.140.13}
\]
by the trace description in A.136. Exactness of the homotopy-fibre
sequence makes this zero on the image of $\pi_1A_{\mathrm{PL}}$.
Equivalently, a jointly PL loop has a jointly PL trace, whose
relative PL-structure difference is zero. Surjectivity of
$\eta_B$ proves the claimed failure of $\pi_1$-surjectivity
whenever the cohomology is nonzero.

For the concrete example $B=S^2\times S^2$, the product cell
decomposition gives $H^2(B;\mathbb Z/2)\cong(\mathbb Z/2)^2$;
it is a closed simply connected finite PL four-manifold. An exact
finite model uses two ordered tetrahedron boundaries. On each
product of two triangular faces, sorting the two pairs of cumulative
barycentric coordinates divides the product into the six simplices
indexed by monotone $(2,2)$-paths. The subdivisions agree on common
ordered faces, so their union is a PL triangulation of the product
itself. Each factor is a PL two-sphere; the product is consequently
a closed PL four-manifold and is simply connected, because the
fundamental group of a product is the product of the fundamental
groups. These are mathematical identifications, not inferences
from a homology calculation.

The exact certificate
\texttt{code/check\_A140\_product\_spheres.py}, with all facets
in \texttt{logs/check\_A140\_product\_spheres.json}, checks the
resulting face vector and mod-$2$ Betti numbers:
\[
 f=(16,84,216,240,96),\qquad b=(1,0,2,0,1).
\]
It also checks boundary-square-zero, all codimension-one incidences,
and the mod-$2$ homology of every nonmaximal-simplex link.
All $1465$ checks passed. The author's independent reading and
rerun agree with the root certificate.
The certificate concerns this finite base and its cohomology,
not an algorithmic construction of the source-theorem TOP loops.
\end{proof}
